\documentclass{beamer}  % Defining this file as a beamer presentation.

\usepackage{url}    % Need this package to display URLs properly.
\usepackage{xcolor}  % Use this package to use colored text.
\usepackage{epsfig} % Use this package to insert figures created externally.

\newcommand{\bi}{\bigskip}  % You can define your own commands to save yourself writing.
\newcommand{\med}{\medskip}  % You can define your own commands to save yourself writing.
\newcommand{\sm}{\smallskip}  % You can define your own commands to save yourself writing.

%\title{An introduction to \LaTeX~for students}
%\author{Christopher Hanusa}
%\date{February 17, 2011}

\setbeamertemplate{theorem begin}  %Need this to 
{%
\begin{\inserttheoremblockenv}
{%
\inserttheoremheadfont
\inserttheoremname
\inserttheoremnumber
\ifx\inserttheoremaddition\@empty\fi%
\inserttheorempunctuation
}%
}
\setbeamertemplate{theorem end}{\end{\inserttheoremblockenv}}


\begin{document}

\frame
{
\centering
{\color{black!30!blue!100} \Large An introduction to \LaTeX\ for students}

\bi
{\large Christopher Hanusa}

\bi
{\large February 17, 2011}

\bi
\url{Christopher.Hanusa@qc.cuny.edu}

\bi
\url{http://people.qc.cuny.edu/chanusa/} $>$ \url{Talks}


}

\section{Introduction}
\subsection{What can you do?}
\frame
{
  \frametitle{Pros and Cons of \LaTeX}

Why use \LaTeX?
\begin{itemize}
\item Ideal for typesetting mathematics.
\item Automatic numbering and citations!
\item Gives the user control of page formatting.
\item Separates the writing of content from the formatting.
\item Free to use; support of the open source community.
\item The output is a pdf; readable by all.
\item It's standard.
\item It's pretty.
\end{itemize}

\bi
When NOT to use \LaTeX:
\begin{itemize}
\item Creating flyers.
\item When you have only text. (No figs, tables, eqns, references)
\item If you don't want to be in charge of the formatting.
\end{itemize}

}


\frame
{
  \frametitle{Typesetting Mathematics}

$$\int_{1}^{9} \frac{\partial}{\partial y}\Big(yx^5+e^{xy}\Big)\, dx$$ 

$$\sum_{n=1}^{\infty} \frac{1}{n^2} =\frac{\pi}{6}$$

$$\sqrt{a^{2} + b^{2}}$$

$$\alpha\beta\gamma\delta\cdots\omega\quad \mathbb{ABC\cdots Z}\quad \mathcal{ABC\cdots Z}$$

$$(A \cap B) \cup C = (A \cup C) \cap (B \cup C)$$

$$\begin{pmatrix} 1 & 2 \\ 3 & 4 \end{pmatrix}\begin{pmatrix} -1 & 0 \\ 3 & 1 \end{pmatrix}=\begin{pmatrix} 5 & 2 \\ 9 & 4 \end{pmatrix}$$

}

\frame
{
\frametitle{Automatic numbering}

I am able to reference theorems, figures, and equations even if their numbers change.



\begin{theorem}[]

I am a theorem.
\label{thm:1}
\end{theorem}
\begin{proof}
Some pretty, distracting equations:
\begin{equation*}\label{eqn0}
|z|=\begin{cases} 
z & \text{ if $z\geq 0$} \\
-z &\text{ if $z<0$}
\end{cases}
\end{equation*}
\begin{equation}\label{eqn1}
z=14x-7y
\end{equation}
\begin{equation}\label{eqn2}
10x=17y
\end{equation}

By Equations~\eqref{eqn1} and \eqref{eqn2}, we have proved Theorem~\ref{thm:1}.
\end{proof}


\vfill
}

\frame
{
\frametitle{Structure of a \url{.tex} file}

\epsfig{figure=texsource, height=3in}

\vspace{-3in}\hspace{2.1in}\begin{minipage}{2.2in}
Always begins with {\color{red} $\backslash$documentclass}

\med
Next, import necessary packages.

\med
Insert user-created commands.

\bi\med
File content starts with

{\color{red} $\backslash$begin\{document\}}


\bi
To break into sections, use 

$\backslash$section and $\backslash$subsection

\bi
Use \% to write comments or to help with the visual structure.


\bi\med
File content ends with

{\color{red} $\backslash$end\{document\} }

\end{minipage}\!\!\!\!

}

\frame
{
\frametitle{Typing text and using fonts}

To type text, type normally;  extra whitespace       does    not 
matter.

To go to a new paragraph, skip a line.  Use a pair of left quotes to open a quote; ``use a pair of right quotes to close.''

\med
The following characters are reserved:
	\#	\$ 	\% 	\& 	\~ 	\_ 	\^{} 	\{ 	\} 	$>$ 	$<$ 	$\backslash$
	
\begin{itemize}
\item $\backslash$ is what tells \LaTeX\ that you are entering a command.
\item \% is for entering comments.
\item \$ is for writing in math mode.
\end{itemize}

\bi
For fonts, surround the text you wish with braces and insert\!\!\!

\sm
{\bf $\backslash$bf} (bold), 
{\em $\backslash$em} (emphasized), 
{\rm $\backslash$rm} (roman), 
{\tt $\backslash$tt} (fixed width), 

\sm
as in {\color{red} \tt \{$\backslash$bf Hello\}} to produce {\bf Hello}.

\vfill
}

\frame
{
\frametitle{Including equations}

\LaTeX\ excels at integrating the writing of mathematics into text.

To include math symbols or equations ``inline'', use {\color{teal} \$} to begin and to end; for example, type {\tt {\color{teal} \$a\_1\^{}2 + b\_1\^{}2\$}} to get $a_1^2+b_1^2$.

\med
Certain symbols everyone uses:

\begin{tabular}{ll}
$\sqrt{10}$ & {\color{teal} \$$\backslash$sqrt\{10\}\$ } \sm\\
$\int_{1}^{10}$ & {\color{teal} \$$\backslash$int\_\{1\}\^{}\{10\}\$ }\sm \\
$\sum_{n=1}^\infty$ & {\color{teal} \$$\backslash$sum\_\{n=1\}\^{}$\backslash$infty\$ }\sm \\
$\frac{1}{n^2}$ & {\color{teal} \$$\backslash$frac\{1\}\{n\^{}2\}\$ }\sm \\
$\alpha$ & {\color{teal} \$$\backslash$alpha\$ } \sm \\
$\hdots$ & {\color{teal} \$$\backslash$hdots\$}; also useful {\color{teal} \$$\backslash$cdots\$} and {\color{teal} \$$\backslash$vdots\$}
\end{tabular}

\med
Each mathematician has their own symbol needs; either peruse tables of symbols or use Detexify 
{\small (\url{detexify.kirelabs.org})}\!\!

\vfill
}

\frame
{
\frametitle{Referencing equations, figures, theorems}

Alternatively, create equations on their own lines using 
{\color{red} $\backslash$begin\{equation\}} and {\color{red} $\backslash$end\{equation\}}, as in 

\sm
{\color{red} $\backslash$begin\{equation\}} {\color{teal} $\backslash$frac\{a\}\{b\}+1=$\backslash$frac\{a+b\}\{b\} }{\color{red} $\backslash$end\{equation\}}

\begin{equation}  \frac{a}{b}+1=\frac{a+b}{b} \end{equation}

\med
Include a {\color{orange} $\backslash$label\{{\em name}\}} to reference it later with {\color{orange} $\backslash$ref\{{\em name}\}}.

\bi
Same goes with defining figures and theorems:

\med
\begin{tabular}{ccc} 
\begin{tabular}{|l|} \hline
{\color{red} $\backslash$begin\{figure\}}\\
~~~~(figure here) \\
~~~~{\color{orange} $\backslash$label\{{\em fig:name1}\}} \\
{\color{red} $\backslash$end\{figure\}}\\ \hline
\end{tabular}
& \qquad\qquad &
\begin{tabular}{|l|}\hline
{\color{red} $\backslash$begin\{theorem\}}\\
~~~~(theorem here) \\
~~~~{\color{orange} $\backslash$label\{{\em thm:name2}\}} \\
{\color{red} $\backslash$end\{theorem\}}\\ \hline
\end{tabular}
\end{tabular}

\bi
$\hdots$ Figure {\color{orange} $\backslash$ref\{{\em fig:name1}\}} exhibits Theorem {\color{orange} $\backslash$ref\{{\em thm:name2}\}} $\hdots$


}



\frame
{
\frametitle{Lists and Tables}

Lists are pretty simple; use {\color{red} itemize} (bullets) or {\color{red} enumerate} (numerals).\!\!\!\!\!\!

\bi
\begin{enumerate}
\item Item 1.
\item Item 2.
\begin{itemize}
\item Item 2a.
\item Item 2b.
\end{itemize}
\end{enumerate}
\vspace{-0.8in}\hspace{1.8in}\begin{minipage}{2in}
\begin{tabular}{|l|} \hline
{\color{red} $\backslash$begin\{enumerate\}}\\
~~~~$\backslash$item  Item 1. \\
~~~~$\backslash$item  Item 2. \\
~~~~{\color{red} $\backslash$begin\{itemize\}}\\
~~~~~~~~$\backslash$item  Item 2a. \\
~~~~~~~~$\backslash$item  Item 2b. \\
~~~~{\color{red} $\backslash$end\{itemize\}}\\ 
{\color{red} $\backslash$end\{enumerate\}}\\ \hline
\end{tabular}
\end{minipage}

\bi
Tables are pretty complicated; use {\color{red} tabular}.

\vspace{0.6in}

}


\frame
{
\frametitle{Miscellaneous}

\begin{itemize}

\item To include graphics, use the package  {\tt \color{red} epsfig} or {\tt \color{red} graphicx}.

\par {\tt $\backslash$epsfig\{figure=residual4.eps,height=1in\}}

\begin{center}
\fbox{\epsfig{figure=residual4.eps, height=1in}}
\end{center}

\item To use color, use the package {\tt \color{red} xcolor}.

For example, \{$\backslash$color\{red\} xcolor\}.

\end{itemize}

\begin{itemize}
\item To organize citations, use Bib\TeX\  
\item To create slide shows, use Beamer
\item LaTeX integrates with the editor \url{emacs}.
\end{itemize}

}

\frame
{
\frametitle{Installation}

So you're hooked!  How to get \LaTeX\ for yourself?

\begin{itemize}
\item  On a Mac
   \begin{itemize}
   \item  Distribution: MacTeX (\url{tug.org/mactex/})
   \item  Viewer: TeXShop (\url{texshop.org})
   \end{itemize}
\item  On a PC
   \begin{itemize}
   \item  Distribution: MiKTeX (\url{miktex.org})
   \item  Viewer: TeXnicCenter (\url{texniccenter.org})
   \end{itemize}
\item  On Linux
   \begin{itemize}
   \item  Like everything else, it's possible.
   \end{itemize}
\end{itemize}

}




\frame
{
\frametitle{Tips and help}

This is one academic exercise where I say: PLAGIARIZE!

Read other people's files to see what they do; then copy.

\med
Search the web!  Example: ``latex tabular''

\med
Helpful reference sheets (condensed):
\begin{itemize}
\item Typing Math: Short Math Guide for \LaTeX\ {\small \url{ftp://ftp.ams.org/pub/tex/doc/amsmath/short-math-guide.pdf}}
\item Find symbols: \url{detexify.kirelabs.org}  or search.
\end{itemize}

\med
Helpful for getting started (some reading):
\begin{itemize}
\item \url{http://www.tug.org/begin.html}
\item \url{www.ctan.org/tex-archive/info/mil/mil.pdf}
\end{itemize}

\vspace{.5in}
}

\frame
{
\frametitle{Let's get hands on!}

\begin{enumerate}
\item Download Empty.tex and NotEmpty.tex from my website.

\item Open each using TeXShop on your computer.

\item Start typing.

\item When you want to see the output, click on the ``Typeset'' button or click Cmd-S (save) Cmd-T (typeset).  

\item This will pop up a small window which runs the ``latex'' command, and (if you have no errors), will pop up a window with the output in pdf format.

\item If you have errors, you need to decode them, fix them, and typeset again.  (often: mismatched \{ \} or misspelled command)\!\!\!\!
\end{enumerate}

\vspace{.5in}

}

\end{document}





































